% Table 2 — Conceptual comparison with the closest prior art.
\begin{table}[t]
\centering
\small
\caption{What the GLR controls compared with the closest prior art. The distinction is the
object (a readout on a frozen force-free propagator) and the target (the canonical mean mode/group
kinetic partition, with a stochastic marginal option, plus the multi-lag velocity ACF), not the
use of a memory kernel per se.}
\label{tab:prior}
\resizebox{\linewidth}{!}{%
\begin{tabular}{@{}p{2.1cm}p{3.1cm}p{4.0cm}p{3.6cm}@{}}
\toprule
Approach & Object & What is controlled or learned & Evidence \\
\midrule
TrajCast & force-free propagator + global CSVR & one-step displacement/velocity; total
temperature & water RDF/VDOS/diffusion \\
Bigi--Ceriotti & deterministic structure-preserving map & symplectic/time-reversible structure;
energy/equipartition & water, GeTe, Ar, Al \\
MAMD & LSTM velocity update & stability via periodic learned velocity injection & several MD
systems \\
GLE thermostat & memory kernel + colored noise in an integrator & FD stationary distribution of
the integrator & generic MD \\
\textbf{GLR} & \textbf{readout on a frozen force-free propagator} & \textbf{mean mode kinetic
partition (+ stochastic marginal) + multi-lag velocity ACF} & \textbf{water NVT, two 2.5 ns
starts; quartz 5 fs, two seeds} \\
\bottomrule
\end{tabular}}
\end{table}
